\section{Finite row sums}\label{supp:rows}
The row sums can be evaluated without listing the $M$ monomials. For example, let
\[
 T(N,m)=\sum_{s=0}^{N-1}(N-s)\floor{s/m},
 \qquad Q=\floor{(N-1)/m}.
\]
Counting each multiple of $m$ once gives the exact formula
\begin{equation}\label{eq:T-floor}
 T(N,m)=\frac12\sum_{a=1}^{Q}(N-am)(N-am+1).
\end{equation}
Thus the tower contribution summed over the triangle is
$K\sum_{a\geq1}T(S,q^{2a})$. The sums terminate when $q^{2a}>S$.
A similar count in the $r$ direction gives, with $Q=\floor{R/q}$,
\begin{equation}\label{eq:Sr}
 S_r(R,K,q)=\sum_{a=1}^{Q}(R-aq)(R+K-aq),
 \qquad
 2\sum_{\Lambda}\floor{r/q}\geq\frac{S_r(R,K,q)}K.
\end{equation}
For each $r$, the number of permitted $s$ is
$S-\lfloor r/K\rfloor=(R-r+(r\bmod K))/K$. Summing this count above
each threshold $r=aq$, in complete blocks of $K$ and one remainder block,
gives the inequality.
The difference between $2K\sum_\Lambda\floor{r/q}$ and $S_r$ is
$\sum_{a=1}^{Q}c_a(K-c_a)$, where $c_a=aq\bmod K$; all these terms are nonnegative. Formula~\eqref{eq:Sr} is a quadratic polynomial in $q$ on each interval where $Q$ is constant. This makes the large prime sums practical to evaluate with integer arithmetic.



\section{Complementary singular products}\label{supp:complementary}
Proof of Lemma~\ref{lem:complementary}.
\begin{proof}
The matrix $T_J$ is triangular with diagonal entries $d^j$. Its full squared singular product is therefore $d^{J(J-1)}$. The largest $m$ squared singular product equals this determinant times the largest $k$ squared singular product of $T_J^{-1}$.

The inverse is composition by $(z-c_0)/d$. Multiplying it on the right by $D_{s_0}=\operatorname{diag}(s_0^j)$ gives composition by
$-c_0r_0/G_0+(r_0/G_0)z$, a strict disk contraction. The negative translation has the same singular
values as its positive counterpart, by the unitary reflection $z\mapsto-z$.
Lemma~\ref{lem:composition} gives squared singular values $G_0^b r_0^{2l}$ for its infinite extension. The norm inequality for the $k$th exterior power of a product bounds the singular product of
$T_J^{-1}=(T_J^{-1}D_{s_0})D_{s_0}^{-1}$ by the product of the two bounds. The largest $k$ entries of $D_{s_0}^{-1}$ occur at indices $J-k,\ldots,J-1$ if $s_0\leq1$ and at $0,\ldots,k-1$ if $s_0\geq1$. Their squared products contribute exactly $-N\log s_0$.
\end{proof}

\section{Parameters for the finite comparisons}\label{supp:parameters}
\begin{table}[ht]
\centering
\begin{tabular}{rrrrr}
\toprule
$P$ & $S$ & $R$ & $J$ & $n$\\
\midrule
$36\,000\,000$ & $1\,004\,651$ & $43\,199\,993$ & $5\,399\,999$ & $4\,822\,400$\\
$169\,000\,000$ & $4\,716\,279$ & $202\,799\,997$ & $25\,349\,999$ & $22\,638\,400$\\
$312\,000\,000$ & $8\,706\,976$ & $374\,399\,968$ & $46\,799\,996$ & $41\,793\,600$\\
$312\,000\,000$ & $8\,706\,976$ & $374\,399\,968$ & $46\,799\,996$ & $35\,394\,400$\\
\bottomrule
\end{tabular}
\caption{Parameters for all determinant comparisons below $875$. The last
line is used on $[499,503]$; the other lines at the individual discriminants
listed in Table~\ref{tab:candidates}.}
\label{tab:parameters}
\end{table}



At $D=499$, the exact difference in \eqref{eq:uniform-rank-check} is
\begin{equation}\label{eq:exact-rank-margin}
 \frac{7\,771\,956\,952\,074\,657\,551}{10\,750\,000}>7.22\times10^{11}.
\end{equation}
